\begin{abstract}
\noindent
This note documents the models used in the paper  \emph{The Environmental Macroeconomics of the Circular Economy}. It includes a planner's problem in continuous time, a corresponding decentralized market economy, and quantitative implementations in discrete time with explicit functional forms. A defining nature of the models are that the material balance principle (mass preservation and increasing entropy) is imposed for each economic activity -- production, extraction, exploration, waste collection and treatment, consumption, investment, and recycling -- with aggregate material balance principles derived from the process-level accounting. Doing so ties the use of materials (through extraction of raw materials from nature and recycled materials from waste) together with the waste generation and pollution and yields an endogenous material intensity of final goods and capital. 


The note contains four parts that deals with the different model types: The planner's problem in continuous time, a corresponding market implementation, the quantitative version of the planner's problem in discrete time with explicit functional forms, and a corresponding market version. 


In the first section, we present the general setup for the theoretical model. In this most general framework, we allow for waste generation to occur whenever final goods are used for an economic activity (representing embodied materials diffused as waste or accumulated in durable stocks that will delay diffusion) \emph{and} proportional to activities measured in tonnes as well (extraction of new resources and exploration for new resources, both in tonnes).

The second part decentralizes that allocation. The market economy fails on three counts rather than one --- the pollution stock is an externality, the material balance is a constraint on the economy that binds no individual firm, and the material intensity of production is an economy-wide object that every agent takes as given --- and we show that three prices are nonetheless enough to repair all three: an emission tax, a charge per tonne of waste released from the capital stock or displaced from nature, and a lower charge per tonne diffused when a final good is used up. The resulting instrument set is a Pigouvian emission tax combined with a deposit--refund scheme on materials embodied in durable capital, and the materials part of it is exactly self-financing.

The third and fourth parts restate the planner's problem and the market economy in discrete time with explicit functional forms, ready for numerical solution. Discretization turns out to require one substantive addition --- because recycled material re-enters production in the period in which the waste was generated, the accounting identities become a simultaneous system inside the period --- and otherwise amounts to replacing each shadow price of a stock by its present value. Under the parameterization adopted there, the recycling technology is a constant-elasticity aggregator of treated waste and recycling capital, the condition for the onset of the circular economy reduces to a comparison between the recycling frontier times the value of a displaced virgin tonne and the marginal product of capital, and the sustainable rate of decline of virgin extraction is bounded by the rate of technical progress in recycling.

The appendix presents a couple of model variants including short solution characterizations that may be either relevant to explore further in the future or be useful for the numerical implementation. We also collect a list of other model extensions/adjustments that may be worth exploring later/in separate projects. 
\end{abstract}